%\clearpage
\section{Límites. Límites superiores e inferiores.}

En esta sección, $(X,d_X)$ y $(Y,d_Y)$ son espacios métricos.

\begin{definicion}[Límite de una función]
Sean $A\subseteq X$ y $f:A\to Y$.
Si $p$ es un punto de acumulación de $A$ y $b\in Y$, decimos que
\[
\lim_{x\to p}f(x)=b
\]
si para todo $\varepsilon>0$ existe $\delta>0$ tal que para todo $x\in A\setminus\{p\}$,
\[
d_X(x,p)<\delta \implies d_Y(f(x),b)<\varepsilon.
\]
\end{definicion}

\begin{remark}
La negación de la definición anterior es:
\[
(\exists\varepsilon>0)(\forall\delta>0)(\exists x_\delta\in A\setminus\{p\})\;
d_X(x_\delta,p)<\delta \ \land\ d_Y(f(x_\delta),b)\ge\varepsilon.
\]
\end{remark}

\begin{teorema}[Caracterización secuencial del límite]
Supongamos que $p$ es punto de acumulación de $A\subseteq X$ y $b\in Y$.
Entonces
\[
\lim_{x\to p}f(x)=b
\]
si y sólo si para toda sucesión $\{x_n\}\subseteq A\setminus\{p\}$ con $x_n\to p$, se tiene $f(x_n)\to b$.
\end{teorema}

\begin{example}
Probar que
\[
\lim_{x\to 2}(3x+1)=7.
\]
Sea $\varepsilon>0$. Si tomamos $\delta=\varepsilon/3$, entonces
\[
|x-2|<\delta \implies |(3x+1)-7|=3|x-2|<\varepsilon.
\]
\end{example}

\subsection*{Límites laterales en funciones reales}

Sea $A\subseteq\R$, $f:A\to\R$ y $p\in A'$.

\begin{definicion}
\begin{enumerate}
    \item Decimos que $\lim_{x\to p^+}f(x)=b$ si para todo $\varepsilon>0$ existe $\delta>0$ tal que, si $x\in A$ y $p<x<p+\delta$, entonces $|f(x)-b|<\varepsilon$.
    \item Decimos que $\lim_{x\to p^-}f(x)=b$ si para todo $\varepsilon>0$ existe $\delta>0$ tal que, si $x\in A$ y $p-\delta<x<p$, entonces $|f(x)-b|<\varepsilon$.
\end{enumerate}
\end{definicion}

\begin{teorema}
Sean $A\subseteq\R$, $p\in A'$, $f:A\to\R$ y $b\in\R$.
Entonces
\[
\lim_{x\to p}f(x)=b
\]
si y sólo si
\[
\lim_{x\to p^+}f(x)=b \quad\text{y}\quad \lim_{x\to p^-}f(x)=b.
\]
\end{teorema}

\begin{teorema}[Teorema de estricción]
Sean $f,g,h$ funciones definidas en $A\subseteq\R$, y sea $p\in A'$.
Supongamos que para $x\in A\setminus\{p\}$ suficientemente cercano a $p$ se cumple
\[
f(x)\le g(x)\le h(x),
\]
y además
\[
\lim_{x\to p}f(x)=L,
\qquad
\lim_{x\to p}h(x)=L.
\]
Entonces
\[
\lim_{x\to p}g(x)=L.
\]
\end{teorema}

\begin{example}
Calcular $\lim_{x\to 0}x^2\sen(1/x)$.

Como $-1\le\sen(1/x)\le 1$ para $x\neq 0$, al multiplicar por $x^2>0$:
\[
-x^2\le x^2\sen(1/x)\le x^2.
\]
Además,
\[
\lim_{x\to 0}(-x^2)=0,
\qquad
\lim_{x\to 0}x^2=0.
\]
Por el teorema de estricción,
\[
\lim_{x\to 0}x^2\sen(1/x)=0.
\]
\end{example}
